% eksperymentalne tłumaczenie na włoski

%

\pol{\section{Zastosowanie: punkt i kąt Crelle'a-Brocarda}}
\ita{\section{Applicazione: punto e angolo di Crelle-Brocard}}

\pol{Poznamy teraz dwa wyróżnione punkty trójkąta, opisane w 1875 roku przez oficera francuskiej armii, Henriego Brocarda oraz wiele lat wcześniej przez Augusta Crelle'a \cite{crelle_1816}, założyciela słynnego czasopisma matematycznego.}
\ita{Studieremo ora due punti notevoli del triangolo, descritti nel 1875 dall'ufficiale dell'esercito francese Henri Brocard e molti anni prima da August Crelle \cite{crelle_1816}, fondatore di una celebre rivista matematica.}
\index[persons]{Brocard, Henri}%
\index[persons]{Crelle, August}%

\begin{definition}
\label{punkty_brocarda}%
    \pol{Niech $\triangle ABC$ będzie trójkątem.}
    \ita{Sia $\triangle ABC$ un triangolo.}
    \pol{Punkt $X$ leżący w jego wnętrzu taki, że kąty $\angle XAB$, $\angle XBC$, $\angle XCA$ są równej miary, nazywamy (pierwszym) punktem Crelle'a-Brocarda.}
    \ita{Il punto $X$ interno al triangolo tale che gli angoli $\angle XAB$, $\angle XBC$, $\angle XCA$ abbiano la stessa ampiezza è chiamato (primo) punto di Crelle-Brocard.}
    \pol{\index{punkt!Crelle'a-Brocarda}}%
    \ita{\index{punto!di Crelle-Brocard}}%
    \begin{center}
\begin{comment}
    \begin{tikzpicture}[scale=.75]
        ...
    \end{tikzpicture}
\end{comment}
    \end{center}
    % Construct a circle through vertices A and B, tangent to side BC of the triangle.
    % Symmetrically, construct the other two circles.
    % These three circles intersect at the first Brocard Point of triangle ABC.
    % https://www.geogebra.org/m/MT679Keu
\end{definition}

\pol{Miarę wspomnianych kątów nazywamy kątem Crelle'a-Brocarda.}
\ita{La misura dei suddetti angoli è chiamata angolo di Crelle-Brocard.}
\pol{\index{kąt!Crelle'a-Brocarda}}%
\ita{\index{angolo!di Crelle-Brocard}}%
\pol{Jest jeszcze drugi punkt Crelle'a-Brocarda, gdzie odwracamy kolejność punktów: $\angle XBA = \angle XCB = \angle XAC$.}
\ita{Esiste anche un secondo punto di Crelle-Brocard, ottenuto invertendo l'ordine dei punti: $\angle XBA = \angle XCB = \angle XAC$.}

\begin{proposition}
\label{miara_kata_crellea_brocarda}%
    \pol{W każdym trójkącie istnieje (jedyny) punkt Crelle'a-Brocarda.}
    \ita{In ogni triangolo esiste un unico punto di Crelle-Brocard.}
    \pol{Miara kąta Crelle'a-Brocarda $\omega$ spełnia związek}
    \ita{La misura dell'angolo di Crelle-Brocard $\omega$ soddisfa la relazione}
    \begin{equation}
        \frac 1 {\tan \omega} = \frac 1 {\tan \alpha} + \frac 1 {\tan \beta} + \frac 1 {\tan \gamma},
    \end{equation}
    \pol{gdzie $\alpha, \beta, \gamma$ to miary kątów trójkąta.}
    \ita{dove $\alpha, \beta, \gamma$ sono le ampiezze degli angoli del triangolo.}
    \pol{Ponadto, $0 \le \omega \le \pi/6$.}
    \ita{Inoltre, $0 \le \omega \le \pi/6$.}
\end{proposition}

\pol{Zetel \cite[s. 176]{zetel_2020} pochwali się, że Szebarszin zauważył, że funkcje trygonometryczne kąta Brocarda można obliczyć prosto i że wynikiem będzie}
\ita{Zetel \cite[p. 176]{zetel_2020} osserva che Šebaršin notò come le funzioni trigonometriche dell'angolo di Brocard possano essere calcolate facilmente e che il risultato è}
\index[persons]{Szebarszin, M N}%
\begin{equation}
    \cos = \frac{a^2 + b^2 + c^2}{2\sqrt{a^2b^2 + a^2c^2 + b^2c^2}}.
\end{equation}

\pol{Na 32. Międzynarodowej Olimpiadzie Matematycznej pojawi się następujące zadanie:}
\ita{Alla 32ª Olimpiade Internazionale della Matematica comparve il seguente problema:}

\begin{problem}
    \pol{Niech $P$ będzie dowolnym punktem wewnątrz trójkąta $\triangle ABC$.}
    \ita{Sia $P$ un punto arbitrario all'interno del triangolo $\triangle ABC$.}
    \pol{Wykazać, że co najmniej jeden z kątów $\angle PAB$, $\angle PBC$, $\angle PCA$ jest mniejszy lub równy $\pi /6$.}
    \ita{Dimostrare che almeno uno degli angoli $\angle PAB$, $\angle PBC$, $\angle PCA$ è minore o uguale a $\pi /6$.}
\end{problem}

\pol{Waldemar Pompe (w $\Delta_{92}^{3}$) przedstawi siedem wskazówek, po których zadanie rozwiąże się nieomal samo.}
\ita{Waldemar Pompe (in $\Delta_{92}^{3}$) presenterà sette suggerimenti dopo i quali il problema si risolverà quasi da solo.}
% https://www.deltami.edu.pl/1992/03/krok-po-kroku/

\pol{Bogdańska, Neugebauer \cite[s. 100]{neugebauer_2018} podadzą jako ćwiczenie, (Zetel \cite[s. 148]{zetel_2020} zaś z dowodem):}
\ita{Bogdańska e Neugebauer \cite[p. 100]{neugebauer_2018} lo propongono come esercizio (mentre Zetel \cite[p. 148]{zetel_2020} ne fornisce una dimostrazione):}

\begin{proposition}
    \pol{Niech $X$ będzie punktem Crelle'a-Brocarda trójkąta $\triangle ABC$.}
    \ita{Sia $X$ il punto di Crelle-Brocard del triangolo $\triangle ABC$.}
    \pol{Niech $R_a$, $R_b$ i $R_c$ oznaczają promienie okręgów opisanych na trójkątach $\triangle XBC$, $\triangle AXC$, $\triangle ABX$, zaś $R$ będzie jak zwykle promieniem okręgu opisanego na trójkącie $\triangle ABC$.}
    \ita{Siano $R_a$, $R_b$ e $R_c$ i raggi delle circonferenze circoscritte ai triangoli $\triangle XBC$, $\triangle AXC$, $\triangle ABX$, mentre $R$ sia il raggio della circonferenza circoscritta al triangolo $\triangle ABC$.}
    \pol{Wtedy}
    \ita{Allora}
    \begin{equation}
        R = \sqrt[3]{R_a R_b R_c}.
    \end{equation}
\end{proposition}

\pol{Peter Yiff \cite{yff_1963} postawi w 1963 roku (!) hipotezę, że $8 \omega^3 \le \alpha \beta \gamma$.}
\ita{Peter Yiff \cite{yff_1963} formulò nel 1963 (!) la congettura che $8 \omega^3 \le \alpha \beta \gamma$.}
\index[persons]{Yiff, Peter}%
\pol{Dowód znajdzie w 1974 roku Faruk Abi-Khuzam \cite{abikhuzam_1974}.}
\ita{La dimostrazione sarà trovata nel 1974 da Faruk Abi-Khuzam \cite{abikhuzam_1974}.}
\index[persons]{Abi-Khuzam, Faruk}%

\pol{Mamy jeszcze mało ciekawy dla kogoś o wiedzy tak nikłej jak my okrąg Crelle'a-Brocarda.}
\ita{Abbiamo inoltre la circonferenza di Crelle-Brocard, poco interessante per chi possiede conoscenze limitate come le nostre.}
% PRZECZYTANO: https://mathworld.wolfram.com/BrocardCircle.html
\pol{\index{okrąg!Crelle'a-Brocarda}}%
\ita{\index{circonferenza!di Crelle-Brocard}}%
\pol{Jego średnicą jest odcinek łączący środek okręgu opisanego z punktem Lemoine'a.}
\ita{Il suo diametro è il segmento che unisce il circocentro al punto di Lemoine.}
\pol{\index{okrąg!opisany}}%
\ita{\index{circonferenza!circoscritta}}%
\pol{\index{punkt!Lemoine'a}}%
\ita{\index{punto!di Lemoine}}%
\pol{Przechodzi przez pierwszy i drugi punkt Crelle'a-Brocarda oraz wierzchołki celowo niezdefiniowanego tu trójkąta Brocarda, co uzasadnia stosowaną czasami nazwę ,,okrąg siedmiu punktów''.}
\ita{Passa per il primo e il secondo punto di Crelle-Brocard e per i vertici del triangolo di Brocard, volutamente non definito qui, il che giustifica il nome talvolta usato di \emph{circonferenza dei sette punti}.}
\pol{\index{okrąg!siedmiu punktów}}%
\ita{\index{circonferenza!dei sette punti}}%
\pol{Ma promień równy}
\ita{Ha raggio pari a}
\begin{equation}
    R \cdot \frac{\sqrt{1 - 4 \sin^2 \omega}}{2 \cos \omega}.
\end{equation}

%